% ===========================================================================
%  user/08-settings/00-chapter.tex
% ===========================================================================
\chapter{Settings and data}\label{cap:settings}

The \voce{Settings} item gathers three tabs: \voce{Config}, \voce{Data} and
\voce{Market rules}. They are the three places where you tell the system
what to work with, and within which limits.

\section{Data: bringing in the prices}\label{sez:data}

Without data, every page of the application is empty. It is the first thing
to do after installing.

\figuraSchermata{1.0}{figure/schermate/data.png}
  {The Data page. On the left the automatic download, on the right the manual
   upload, below them the categories.}
  {data}

\subsection{Downloading}

You choose three things and press \pulsante{Download}:

\begin{description}[leftmargin = 1.4em, style = nextline, itemsep = 4pt]
  \item[\campo{Provider}] Where to get them from. Some need a personal key,
    which you set up from the site.
  \item[\campo{Timeframe}] The time step. The same instrument can exist at
    several steps: daily data and five-minute data live in different folders
    and do not get in each other's way.
  \item[\campo{Category}] What kind of data it is.
\end{description}

\noindent
The job goes into the queue and the bell tells you when it is done. Below the
fields, a line says what is already there: how many files, how much space, up
to when they go.

\subsection{Uploading your own file}

The \voce{Your data} panel accepts CSV, Parquet and XLSX files dragged onto
it. The system recognises by itself columns written in capitals, dates given
as an index or as a column, and prefixes such as \campo{adj\_}. The
\voce{Expected format} panel shows what it expects.

\subsection{The categories}

Each category is a folder, and they come in two types.

\begin{center}
\begin{tabular}{@{}L{22mm}L{116mm}@{}}
\toprule
\textbf{Type} & \textbf{It means} \\
\midrule
\rigaTenue
\voce{1:1} & One file per instrument, named after the instrument. The link
  is automatic, and the application can say how many instruments are covered
  and how many are not. \\
\voce{1:N} & Files shared by several instruments: sectors, macroeconomic
  series, currencies. Here there is no coverage to compute, and the strategy
  has to declare which file to use. \\
\bottomrule
\end{tabular}
\end{center}

\begin{nota}{\enquote{0 / 15 assets covered} is not a fault}
It means the category exists but has no data yet for any instrument. It
fills up by downloading, or by dragging files into it.
\end{nota}

\section{Config: the limits}\label{sez:config}

These are the parameters that apply to every calculation.

\figuraSchermata{1.0}{figure/schermate/config.png}
  {The parameters. Each field carries its unit below it, because the
   difference between \num{0.12} and \num{12} is a factor of a hundred.}
  {config}

\begin{center}
\begin{tabular}{@{}L{44mm}L{94mm}@{}}
\toprule
\textbf{Field} & \textbf{What it limits} \\
\midrule
\multicolumn{2}{@{}l}{\textit{\small Portfolio \& allocator}}\\[2pt]
\rigaTenue
\voce{Initial capital} & The capital the calculations start from. \\
\voce{Annual target volatility} & How much you accept the portfolio swinging
  in a year. It is the main risk dial. \\
\rigaTenue
\voce{Max weight per position} & The most a single instrument can weigh. \\
\voce{Max gross per asset class} & The most a whole class can weigh ---
  shares, bonds, currencies. It guards against diversification that is only
  apparent: four bonds all going the same way are a single position. \\
\rigaTenue
\voce{Max gross leverage} & The total exposure. \num{1.0} means no
  leverage. \\

\addlinespace
\multicolumn{2}{@{}l}{\textit{\small Ensemble}}\\[2pt]
\rigaTenue
\voce{Target volatility} & The volatility every stream is brought to
  \emph{before} being weighted. \\
\voce{Regime window}, \voce{thresholds} & How the state of the market is
  assessed for the modulation. \\
\bottomrule
\end{tabular}
\end{center}

\begin{attenzione}{\voce{Target volatility} also acts on the real account}
It is not a display setting. The same factor that normalises the streams in
the backtest multiplies the positions the runner opens on a real account.
Changing it changes how much you really risk, not just what you see on the
chart.
\end{attenzione}

\noindent
Further down, \voce{Asset categories} lets you add categories of
instruments: once created, a strategy can declare them and restrict itself to
them.

\section{Market rules: the regimes}\label{sez:regimes}

A \voce{regime} is the state of the market, with three values: \voce{bull},
\voce{neutral}, \voce{bear}. Strategies can declare that they weigh things
differently according to the state.

\figuraSchermata{1.0}{figure/schermate/regime.png}
  {The configured markets. Below them, the history of the states: green bull,
   grey neutral, red bear.}
  {regime}

\subsection{How to read the page}

Each panel at the top is a \emph{market} --- a set of instruments that share
the same way of measuring the regime:

\begin{description}[leftmargin = 1.4em, style = nextline, itemsep = 4pt]
  \item[\voce{RULE}] The rule used to compute the state.
  \item[\voce{DRIVER}] The series that drives the decision --- for shares,
    SPY; for forex, the dollar.
  \item[\voce{SYMBOLS}] How many name patterns fall into this market.
  \item[The coloured bar] How much time was spent in each state.
\end{description}

\begin{attenzione}{\enquote{No driver available}}
The panel of a market that cannot find its driving series says so with a
warning, like the \voce{crypto} market in the screenshot. That market
\textbf{modulates nothing}: it is not an error that blocks anything, but the
strategies on those instruments behave as if the regime were always neutral.
The fix is to download the missing series from the \voce{Data} page.
\end{attenzione}

\subsection{The chart}

It shows the driving series with the background coloured by state. Below it,
the counts: how many observations, how many times the state changed, and the
percentage of time spent in each.

\begin{nota}{Few changes of state is a good thing}
A regime must be a \emph{state}, not a reaction. Twenty-five transitions in
five years mean the rule tells different periods apart; two hundred would
mean it is chasing noise, and a strategy that changes weight every week pays
costs without taking a position.
\end{nota}
